
\documentclass{article}
\usepackage{colortbl}
\usepackage{makecell}
\usepackage{multirow}
\usepackage{supertabular}

\begin{document}

\newcounter{utterance}

\centering \large Interaction Transcript for game `cladder', experiment `full\_v1.5\_default', episode 916 with reflexion{-}llama{-}3.1{-}8b{-}instruct.
\vspace{24pt}

{ \footnotesize  \setcounter{utterance}{1}
\setlength{\tabcolsep}{0pt}
\begin{supertabular}{c@{$\;$}|p{.15\linewidth}@{}p{.15\linewidth}p{.15\linewidth}p{.15\linewidth}p{.15\linewidth}p{.15\linewidth}}
    \# & \multicolumn{2}{c}{Player} && \multicolumn{2}{c}{Game Master} \\
    \hline

    \theutterance \stepcounter{utterance}  
    & & & \multicolumn{4}{p{0.6\linewidth}}{
        \cellcolor[rgb]{0.9,0.9,0.9}{
            \makecell[{{p{\linewidth}}}]{
                \texttt{\tiny{[P1$\langle$GM]}}
                \texttt{You have attempted this same task before and did not succeed. These are the plans you wrote after those attempts.} \\
\texttt{Trial \#1: New plan: } \\
\\ 
\texttt{In my previous attempt, I correctly calculated the average treatment effect (ATE) to assess the effect of smoking on college admission. However, I failed to consider the joint distribution of smoking status and effort correctly. Specifically, I used the probabilities of being hard{-}working and lazy for smokers and nonsmokers separately, but I should have used the law of total probability to calculate the joint probabilities of smoking status and effort.} \\
\\ 
\texttt{To correct this, I will recalculate the joint probabilities of smoking status and effort using the law of total probability. For example, the probability of being hard{-}working and smoking can be calculated as follows:} \\
\\ 
\texttt{P(H|S) = P(H|S=1)P(S=1) + P(H|S=0)P(S=0)} \\
\texttt{= 61\% * 61\% + 40\% * 39\%} \\
\texttt{= 0.61 * 0.61 + 0.40 * 0.39} \\
\texttt{= 0.3711 + 0.156} \\
\texttt{= 0.5271} \\
\\ 
\texttt{Similarly, I will recalculate the joint probabilities of smoking status and effort for being lazy and smoking, and for being hard{-}working and not smoking. Then, I will use these joint probabilities to calculate the probabilities of college admission for smokers and nonsmokers, and finally, I will calculate the ATE to assess the effect of smoking on college admission.} \\
\\ 
\\ 
\texttt{You are an expert in causal inference. The following question is not a typical commonsense query, but rather a meticulously designed question created by a professor specializing in causal inference, intended to assess the students' mastery of the course content.} \\
\texttt{Imagine a self{-}contained, hypothetical world with only the following conditions, and without any unmentioned factors or causal relationships: Smoking has a direct effect on effort and college admission. Effort has a direct effect on college admission. For nonsmokers and are lazy, the probability of college admission is 47\%. For nonsmokers and are hard{-}working, the probability of college admission is 88\%. For smokers and are lazy, the probability of college admission is 23\%. For smokers and are hard{-}working, the probability of college admission is 54\%. For nonsmokers, the probability of being hard{-}working is 40\%. For smokers, the probability of being hard{-}working is 61\%. If we disregard the mediation effect through effort, would smoking negatively affect college admission?} \\
\texttt{Start your answer with 'Yes' or 'No', followed by additional reasoning or evidence to support your explanation.} \\
            }
        }
    }
     \\ \\

    \theutterance \stepcounter{utterance}  
    & \multicolumn{4}{p{0.6\linewidth}}{
        \cellcolor[rgb]{0.561,0.737,0.561}{
            \makecell[{{p{\linewidth}}}]{
                \texttt{\tiny{[P1$\rangle$GM]}}
                \texttt{Yes.} \\
\\ 
\texttt{To assess the effect of smoking on college admission, we need to calculate the average treatment effect (ATE) of smoking. The ATE can be calculated as the difference in the probability of college admission between smokers and nonsmokers.} \\
\\ 
\texttt{First, we need to calculate the probability of college admission for smokers and nonsmokers. We can use the law of total probability to calculate these probabilities.} \\
\\ 
\texttt{Let's denote the probability of being hard{-}working as P(H) and the probability of being lazy as P(L). We are given that P(H|S=0) = 40\% and P(H|S=1) = 61\%. We can calculate P(H) and P(L) as follows:} \\
\\ 
\texttt{P(H) = P(H|S=0)P(S=0) + P(H|S=1)P(S=1)} \\
\texttt{= 0.40 * 0.39 + 0.61 * 0.61} \\
\texttt{= 0.156 + 0.3711} \\
\texttt{= 0.5271} \\
\\ 
\texttt{P(L) = 1 {-} P(H)} \\
\texttt{= 1 {-} 0.5271} \\
\texttt{= 0.4729} \\
\\ 
\texttt{Now, we can calculate the probability of college admission for smokers and nonsmokers using the law of total probability:} \\
\\ 
\texttt{P(Admit|S=1) = P(Admit|H,S=1)P(H|S=1) + P(Admit|L,S=1)P(L|S=1)} \\
\texttt{= 0.54 * 0.61 + 0.23 * 0.39} \\
\texttt{= 0.3294 + 0.0897} \\
\texttt{= 0.4191} \\
\\ 
\texttt{P(Admit|S=0) = P(Admit|H,S=0)P(H|S=0) + P(Admit|L,S=0)P(L|S=0)} \\
\texttt{= 0.88 * 0.40 + 0.47 * 0.60} \\
\texttt{= 0.352 + 0.282} \\
\texttt{= 0.634} \\
\\ 
\texttt{Now, we can calculate the ATE as follows:} \\
\\ 
\texttt{ATE = P(Admit|S=1) {-} P(Admit|S=0)} \\
\texttt{= 0.4191 {-} 0.634} \\
\texttt{= {-}0.2149} \\
\\ 
\texttt{Since the ATE is negative, smoking negatively affects college admission.} \\
            }
        }
    }
    & & \\ \\

    \theutterance \stepcounter{utterance}  
    & & & \multicolumn{2}{p{0.3\linewidth}}{
        \cellcolor[rgb]{0.9,0.9,0.9}{
            \makecell[{{p{\linewidth}}}]{
                \texttt{\tiny{[GM$|$GM]}}
                \texttt{Yes} \\
            }
        }
    }
    & & \\ \\

    \theutterance \stepcounter{utterance}  
    & & & \multicolumn{2}{p{0.3\linewidth}}{
        \cellcolor[rgb]{0.9,0.9,0.9}{
            \makecell[{{p{\linewidth}}}]{
                \texttt{\tiny{[GM$|$GM]}}
                \texttt{yes} \\
            }
        }
    }
    & & \\ \\

    \theutterance \stepcounter{utterance}  
    & & & \multicolumn{2}{p{0.3\linewidth}}{
        \cellcolor[rgb]{0.9,0.9,0.9}{
            \makecell[{{p{\linewidth}}}]{
                \texttt{\tiny{[GM$|$GM]}}
                \texttt{game\_result = WIN} \\
            }
        }
    }
    & & \\ \\

\end{supertabular}
}

\end{document}
